% PLANTILLA MAESTRA — INVESTIGACIÓN APLICADA A LA CONTADURÍA / UAS
% Duplique este archivo por actividad y sustituya todos los campos editables.
% No invente datos, resultados, normas, fuentes ni metadatos curriculares.
\documentclass[spanish,letterpaper,oneside]{article}

% DATOS EDITABLES
\def\documenttitle {Título del proyecto o actividad}
\def\documentsubtitle {Investigación aplicada a la contaduría}
\def\documentsubject {Actividad [POR COMPLETAR]}
\def\documentauthor {Nombre del estudiante}
\def\coursename {Investigación aplicada a la contaduría}
\def\coursecode {Clave por confirmar}
\def\universityname {Universidad Autónoma de Sinaloa}
\def\universityfaculty {Licenciatura en Contaduría}
\def\universitydepartment {Unidad académica por confirmar}
\def\universitydepartmentimage {UAS/assets/marca-uas-provisional}
\def\universitydepartmentimagecfg {height=1.45cm}
\def\universitylocation {Sinaloa, México}
\def\authortable {
  \begin{tabular}{ll}
    \textbf{Estudiante:} & Nombre por completar \\
    Matrícula: & Por completar \\
    \textit{Docente:} & Nombre por completar \\
    Grupo y periodo: & Por completar \\
    Fecha de entrega: & \today \\
    Lugar: & \universitylocation
  \end{tabular}
}

\input{template}
\usepackage{helvet}
\renewcommand{\familydefault}{\sfdefault}
\usepackage{tikz}
\usetikzlibrary{arrows.meta,positioning,shapes.geometric}
\setcitestyle{authoryear,open={(},close={)}}

\definecolor{uasBlue}{HTML}{123B6D}
\definecolor{uasGold}{HTML}{C79A2B}
\newcommand{\porcompletar}[1]{\textcolor{uasBlue}{\textbf{[POR COMPLETAR: #1]}}}

\begin{document}
\templatePortrait
\templatePagecfg
\onehalfspacing

\begin{abstractd}
  \textbf{Problema:} \porcompletar{delimitar el problema contable, su contexto y la unidad de análisis.}

  \textbf{Objetivo:} \porcompletar{declarar el resultado verificable que persigue la investigación.}

  \textbf{Método:} \porcompletar{indicar enfoque, diseño, fuentes o participantes y técnica de análisis.}

  \textbf{Resultado o aporte:} \porcompletar{resumir el hallazgo real o, si es protocolo, el aporte esperado sin presentarlo como resultado.}
\end{abstractd}

\templateIndex
\templateFinalcfg

\section{Introducción}
\porcompletar{presentar el contexto, la relevancia contable, el problema, el objetivo y la ruta del documento. Evitar iniciar con definiciones genéricas.}

La investigación aplicada exige que las decisiones metodológicas respondan al problema y que las conclusiones se limiten a la evidencia disponible \citep{hernandezSampieriMetodologia2018,creswellResearchDesign2018}.

\section{Planteamiento del problema}
\subsection{Contexto y antecedentes}
\porcompletar{identificar organización, sector, periodo, antecedentes y evidencia inicial verificable.}

\subsection{Delimitación y pregunta de investigación}
\porcompletar{precisar alcance temático, espacial, temporal y poblacional. Formular una pregunta investigable.}

\subsection{Justificación}
\porcompletar{explicar utilidad académica, profesional, organizacional o social; no confundir relevancia con promesas no demostrables.}

\section{Objetivos}
\subsection{Objetivo general}
\porcompletar{usar un verbo observable y mantener correspondencia directa con la pregunta.}

\subsection{Objetivos específicos}
\begin{enumerate}
  \item \porcompletar{primer objetivo específico.}
  \item \porcompletar{segundo objetivo específico.}
  \item \porcompletar{tercer objetivo específico.}
\end{enumerate}

\section{Marco de referencia}
\subsection{Antecedentes de investigación}
\porcompletar{comparar estudios previos por problema, método, resultados y vacío que permanece.}

\subsection{Marco conceptual y normativo}
\porcompletar{definir categorías contables y explicar la función de cada norma o criterio aplicable. Toda afirmación técnica requiere fuente.}

\section{Diseño metodológico}
\subsection{Enfoque y diseño}
\porcompletar{justificar enfoque cuantitativo, cualitativo o mixto; alcance y diseño.}

\subsection{Unidad de análisis, población y muestra}
\porcompletar{definir criterios de inclusión, selección y tamaño. No declarar representatividad sin sustento.}

\subsection{Técnicas, instrumentos y procedimiento}
\porcompletar{describir recolección, validación, aplicación, resguardo y análisis de datos.}

\subsection{Ética y confidencialidad}
\porcompletar{explicar consentimiento, anonimización, permisos, manejo de conflictos de interés y protección de información contable.}

\section{Matriz de congruencia}
\begingroup
\small
\setlength{\tabcolsep}{4pt}
\renewcommand{\arraystretch}{1.25}
\begin{longtable}{@{}p{0.18\linewidth}p{0.18\linewidth}p{0.18\linewidth}p{0.18\linewidth}p{0.18\linewidth}@{}}
\caption{Correspondencia metodológica del proyecto}\label{tab:congruencia}\\
\toprule
\textbf{Pregunta} & \textbf{Objetivo} & \textbf{Categorías o variables} & \textbf{Técnica} & \textbf{Evidencia esperada} \\
\midrule
\endfirsthead
\toprule
\textbf{Pregunta} & \textbf{Objetivo} & \textbf{Categorías o variables} & \textbf{Técnica} & \textbf{Evidencia esperada} \\
\midrule
\endhead
\bottomrule
\endlastfoot
\porcompletar{pregunta} & \porcompletar{objetivo} & \porcompletar{categorías} & \porcompletar{técnica} & \porcompletar{evidencia} \\
\end{longtable}
\endgroup

\section{Resultados o resultados esperados}
\porcompletar{si ya existe trabajo de campo, presentar datos reales y su fuente; si es protocolo, identificar claramente resultados esperados.}

\section{Discusión}
\porcompletar{interpretar los resultados, contrastarlos con antecedentes y explicar implicaciones para la práctica contable. Separar descripción de interpretación.}

\section{Conclusiones y recomendaciones}
\porcompletar{responder la pregunta y cada objetivo con base en la evidencia, declarar limitaciones y proponer acciones viables. No introducir datos nuevos.}

\section*{Lista de control antes de entregar}
\begin{itemize}
  \item La pregunta, los objetivos, el método, la evidencia y las conclusiones son congruentes.
  \item Los datos son reales, trazables y se distingue entre resultado observado y expectativa.
  \item Se protege la información confidencial y se declaran permisos y limitaciones.
  \item Cada tabla o figura tiene título, fuente y explicación.
  \item Toda cita aparece en la bibliografía y no hay referencias decorativas.
  \item Se eliminaron todos los marcadores \texttt{[POR COMPLETAR]}.
\end{itemize}

\bibliography{investigacion-aplicada-a-la-contaduria}
\end{document}
